\clearpage
\section{Prime-Packet Differentials and Radical Divisibility}

\subsection{The prime-coordinate module}
Let
\[
 a+b=c,
 \qquad
 \gcd(a,b)=1.
\]
The three integers are pairwise coprime. Write
\[
 S=S_a\sqcup S_b\sqcup S_c
 =\{p:p\mid abc\}.
\]
The first object of the calculus is the free prime-packet module
\[
 \Omega_S=\bigoplus_{p\in S}\Z\xi_p.
\]
First define the global free module over all primes, and restrict it to $\Omega_S$ for $S$-supported integers. For such a positive integer $n$, define the universal multiplicative differential
\[
 \boxed{
 \mathrm dn
 =
 n\sum_{p\mid n}\frac{v_p(n)}p\,\xi_p.
 }
\]
Every coefficient is an integer because $p\mid n$. The map obeys the exact Leibniz law
\[
 \mathrm d(mn)=m\,\mathrm dn+n\,\mathrm dm.
\]
The verification is local. For each prime $p$,
\[
 v_p(mn)=v_p(m)+v_p(n),
\]
so the $p$-coordinate of the right side is
\[
 mn\frac{v_p(m)+v_p(n)}p.
\]

A scalar packet is an integer vector
\[
 x=(x_p)_{p\in S}\in\Z^S.
\]
Pairing with $x$ gives the arithmetic derivative
\[
 \boxed{
 \mathrm d_x(n)
 =
 n\sum_{p\mid n}\frac{v_p(n)}p x_p.
 }
\]
The Leibniz rule survives scalarization.

Arithmetic derivatives and their connection with $abc$ have an established literature. The present construction is organized here in prime-packet coordinates and later refined into filtered shell objects~\cite{PastenDerivatives,VojtaABC}.

\subsection{The additive row}
The condition
\[
 \mathrm d_x(a)+\mathrm d_x(b)=\mathrm d_x(c)
\]
is one integer equation. Define
\[
 \lambda_p=
 \begin{cases}
 \dfrac{a v_p(a)}p,&p\mid a,\\[6pt]
 \dfrac{b v_p(b)}p,&p\mid b,\\[6pt]
 -\dfrac{c v_p(c)}p,&p\mid c.
 \end{cases}
\]
Then
\[
 \boxed{\lambda\cdot x=0.}
\]
The additive derivative lattice is
\[
 \boxed{
 \mathcal T_{a,b,c}
 =
 \ker(\lambda)\cap\Z^S.
 }
\]
Since the row is nonzero, its rank is $|S|-1$. The problem is not existence of additive derivatives. The problem is finding an additive derivative that remains transverse to the dependence direction introduced by $a$ and $b$.

\subsection{The arithmetic Wronskian}
Define
\[
 \boxed{
 W_x(a,b)
 =a\,\mathrm d_x(b)-b\,\mathrm d_x(a).
 }
\]
Writing
\[
 \mu_p=
 \begin{cases}
 -\dfrac{v_p(a)}p,&p\mid a,\\[6pt]
 \dfrac{v_p(b)}p,&p\mid b,\\[6pt]
 0,&p\mid c,
 \end{cases}
\]
we obtain
\[
 W_x(a,b)=ab\,\mu\cdot x.
\]
An additive packet is \emph{transverse} when $W_x(a,b)\neq0$. It is degenerate when the two input terms remain derivative-proportional.

Over $\R$, the additive projection is
\[
 P_{\mathrm{add}}
 =
 I-
 \frac{\lambda\lambda^{\mathsf T}}{\norm{\lambda}_2^2}.
\]
The obstruction direction inside the additive hyperplane is
\[
 \mu_\perp=P_{\mathrm{add}}\mu.
\]
When $\mu_\perp\neq0$, the rank-one obstruction projector is
\[
 P_{\mathrm{obs}}
 =
 \frac{\mu_\perp\mu_\perp^{\mathsf T}}
 {\norm{\mu_\perp}_2^2}.
\]
This is the first appearance of the projector-resolved pattern. A large ambient prime space is decomposed into one equation direction, one obstruction direction, and the remaining null packets.

\subsection{Radical divisibility}
Set
\[
 R=\rad(abc).
\]
The central first-order theorem is an exact divisibility law.

\begin{theorem}[Wronskian divisibility]
For every additive packet $x\in\mathcal T_{a,b,c}$,
\[
 \boxed{
 \frac{abc}{R}\mid W_x(a,b).
 }
\]
\end{theorem}

\begin{proof}
Let $p^e\parallel a$. Every term of $\mathrm d_x(a)$ is divisible by $p^{e-1}$. Since $p\nmid b$, the term $b\,\mathrm d_x(a)$ is divisible by $p^{e-1}$. Also $p^e\mid a$, so $a\,\mathrm d_x(b)$ is divisible by $p^e$, hence by $p^{e-1}$. Therefore
\[
 p^{e-1}\mid W_x(a,b).
\]
The same argument applies to primes dividing $b$.

Now let $p^e\parallel c$. Additivity gives
\[
 \mathrm d_x(c)=\mathrm d_x(a)+\mathrm d_x(b).
\]
Using $c=a+b$,
\[
 \begin{aligned}
 W_x(a,b)
 &=a\,\mathrm d_x(b)-b\,\mathrm d_x(a)\\
 &=a\bigl(\mathrm d_x(c)-\mathrm d_x(a)\bigr)-b\,\mathrm d_x(a)\\
 &=a\,\mathrm d_x(c)-c\,\mathrm d_x(a).
 \end{aligned}
\]
The first term is divisible by $p^{e-1}$ because $\mathrm d_x(c)$ is. The second is divisible by $p^e$. Hence $p^{e-1}\mid W_x(a,b)$.

The prime supports of $a$, $b$, and $c$ are disjoint. Multiplying all local divisibilities gives
\[
 \prod_{p^e\parallel abc}p^{e-1}
 =
 \frac{abc}{R}
 \mid W_x(a,b).
\]
\end{proof}

For a transverse packet,
\[
 \boxed{
 |W_x(a,b)|\ge\frac{abc}{R}.
 }
\]
This is the integer analogue of the repeated-root divisor appearing in the polynomial Mason--Stothers theorem.

\subsection{The master height inequality}
For every positive $n$,
\[
 \begin{aligned}
 |\mathrm d_x(n)|
 &\le
 n\norm{x}_\infty
 \sum_{p\mid n}\frac{v_p(n)}p.
 \end{aligned}
\]
Since $p\ge2$,
\[
 \frac1p\le\frac1{2\log2}\log p.
\]
Therefore
\[
 \sum_{p\mid n}\frac{v_p(n)}p
 \le
 \frac{\log n}{2\log2},
\]
and
\[
 |\mathrm d_x(n)|
 \le
 \frac{n\log n}{2\log2}\norm{x}_\infty.
\]
Consequently,
\[
 \begin{aligned}
 |W_x(a,b)|
 &\le
 a|\mathrm d_x(b)|+b|\mathrm d_x(a)|\\
 &\le
 \frac{ab}{\log2}\norm{x}_\infty\log c.
 \end{aligned}
\]
Combining this with radical divisibility gives
\[
 \boxed{
 c
 \le
 \frac{R}{\log2}\norm{x}_\infty\log c.
 }
\]

Except for $(1,1,2)$, define the shortest transverse packet
\[
 \boxed{
 \Theta(a,b,c)
 =
 \min\left\{
 \norm{x}_\infty:
 x\in\mathcal T_{a,b,c},
 W_x(a,b)\neq0
 \right\}.
 }
\]
Then
\[
 c
 \le
 \frac{R}{\log2}\Theta(a,b,c)\log c.
\]
The exceptional triple $(1,1,2)$ has zero input derivatives and is omitted from $\Theta$. It is absorbed into the constant in the conjecture. A subpolynomial estimate
\[
 \Theta(a,b,c)=c^{o(1)}
\]
would imply the strong $abc$ inequality. Later sections explain why this target is too close to the conjecture to be reached by generic lattice geometry alone.

\subsection{Three-sector compression}
For the empty unit sector, set $g_1=0$, $t_1=0$, and $\rho_1(0)=0$. Any nonzero target in that sector is infeasible. Within a nonempty $a$-sector, define the integer row
\[
 \lambda^{(a)}
 =
 \left(\frac{a v_p(a)}p\right)_{p\mid a}
\]
and let
\[
 g_a=\gcd\left\{\frac{a v_p(a)}p:p\mid a\right\}.
\]
Define $g_b$ and $g_c$ similarly. A one-row Smith normal form in each sector produces a unimodular coordinate system in which only one aggregate channel survives. The additive equation becomes
\[
 \boxed{
 g_at_a+g_bt_b-g_ct_c=0.
 }
\]
The Wronskian becomes
\[
 \boxed{
 a g_bt_b-b g_at_a\neq0.
 }
\]
All remaining sector coordinates are lift gauges. They matter only because they determine the cost of realizing an aggregate channel by an integral prime packet.

For each term define
\[
 \rho_a(t)
 =
 \min\left\{
 \norm{x^{(a)}}_\infty:
 \lambda^{(a)}\cdot x^{(a)}=g_at
 \right\},
\]
with analogous functions for $b$ and $c$. Then
\[
 \Theta(a,b,c)
 =
 \min
 \max\{\rho_a(t_a),\rho_b(t_b),\rho_c(t_c)\}
\]
over aggregate solutions satisfying additivity and transversality.

This reduction separates the global two-dimensional arithmetic geometry from the local lifting cost. It is the template used repeatedly later: compress first, then retain the exact synthesis norm of the discarded coordinates.

\subsection{Worked example: \texorpdfstring{$3+125=128$}{3+125=128}}
The prime packets are $3$, $5$, and $2$. The additive row is
\[
 \lambda=(1,75,-448).
\]
A shortest transverse packet is
\[
 x=(-2,6,1).
\]
It gives
\[
 \mathrm d_x(3)=-2,
 \qquad
 \mathrm d_x(125)=450,
 \qquad
 \mathrm d_x(128)=448.
\]
Thus additivity holds. The Wronskian is
\[
 W_x(3,125)
 =3(450)-125(-2)
 =1600.
\]
The radical is
\[
 R=2\cdot3\cdot5=30,
\]
and
\[
 \frac{3\cdot125\cdot128}{30}=1600.
\]
Hence the radical lower bound is attained exactly.

The minimality is elementary. At height at most five, $|x_3+75x_5|\le380<448$, so additivity forces $x_2=0$. Then $x_3=-75x_5$ and the same height restriction force both coordinates to vanish. Thus no transverse packet has height below six, and the displayed packet proves $\Theta=6$.

This example becomes a recurring calibration point. Every later refinement must explain why the sharp integer divisor is visible while the corresponding height bound remains nontrivial.

